\section*{Abstract}

As interest in formalized programming language semantics increases,
proof assistants like Lean offer promising avenues for computer science education.
While other studies frequently focuses on developing new pedagogical tools,
this thesis investigates the structural alignment between an established semantics
textbook and our mechanized translation constructed
to mimic the text as closely as possible.
The study evaluates whether traditional pen-and-paper proofs transfer seamlessly
into Lean or require significant modifications to account for implicit assumptions.
The results indicate that a large portion of the chapters covered in the textbooks
material maps directly into Lean.
However, key deviations emerged due to textbook imprecisions,
omitted "Straightforward" steps, and Lean's specific constraints regarding partial functions.
Ultimately, this close structural mapping shows that mechanized translations
can successfully augment traditional coursework,
provided students receive adequate guidance to manage the rigorous demands
of tool-based formalization.